% ============================================================
%  第四章
% ============================================================
\section{超对称、中心荷，与一个错误的结论}
\markboth{第四章\ 超对称、中心荷，与一个错误的结论}{第四章\ 超对称、中心荷，与一个错误的结论}

\noindent{\small\color{shuyiGray}对应底本：原 p.19--p.24。}\par\medskip

\srcp{19}我在 GNO 论文里感到不安的那一点，直到将近三十年后才被解决；但 Olive 和我实际上成功解决了我对 Montonen--Olive 论文的技术异议。他们的技术假定，在他们所假定的普通（或“玻色”）规范理论里并不成立，但在该理论的超对称版本里确实成立。粗略地说，超对称理论是函数空间上\emph{微分形式}的理论，而不是函数空间上\emph{函数}的理论。而我们证明了：只要有至少 $\cN=2$ 的超对称，Montonen 与 Olive 的经典质量公式在量子力学意义上就真的成立。

（四维时空的超对称规范理论按超对称的量分类：$\cN=1$、$2$ 或 $4$。$\cN=4$ 理论对给定的规范群是唯一的；$\cN=2$ 理论并不完全由规范群决定，但受到高度限制。$\cN=1$ 理论受的限制更少，也更接近现实物理。极粗略地说，$\cN=1,2,4$ 分别对应复数、四元数、八元数。）

\subsection{那篇论文}

\srcp{20}由此产生的那篇论文，大概是到那时为止我写过的最重要的一篇。

\begin{abstractbox}{Witten \& Olive, \emph{Supersymmetry Algebras That Include Topological Charges}（Physics Letters B78, 1978, pp.~97--99）}
We show that in supersymmetric theories with solitons, the usual supersymmetry algebra is not valid; the algebra is modified to include the topological quantum numbers as central charges. Using the corrected algebra, we are able to show that in certain four dimensional gauge theories, there are no quantum corrections to the classical mass spectrum. These are theories for which Bogomolny has derived a classical bound; the argument involves showing that Bogomolny's bound is valid quantum mechanically and that it is saturated.

\medskip
我们证明：在有孤子的超对称理论里，通常的超对称代数不成立；该代数要被修正，以把拓扑量子数作为\emph{中心荷}包含进来。利用修正后的代数，我们可以证明：在某些四维规范理论里，经典质量谱没有量子修正。这些理论正是 Bogomolny 为之导出经典下界的理论；论证的关键在于证明 Bogomolny 下界在量子力学意义上成立，并且被饱和。
\end{abstractbox}

\noindent 但我从中得出了错误的结论。我以为，我们不靠他们那个激进的对偶假定，就解释了 Montonen 与 Olive 的公式，因此并没有真正的证据支持对偶。

\srcp{21}即使我当时得出了正确的结论，我也不认为理解 Montonen--Olive 对偶的时机已经成熟，所以我怀疑那段时间我也做不出多大进展。但在那些早年里，还是做了几件值得注意的事。

\subsection{只有 $\cN=4$ 才对}

\srcp{22}其中之一是：Hugh Osborn 更仔细地考察了单极子的量子化，发现只有 $\cN=4$ 理论（而不是 Olive 和我考虑过的更一般的 $\cN=2$ 理论）才具有 Montonen--Olive 对偶所需的那种性质。

\begin{abstractbox}{Osborn, \emph{Topological Charges for $\cN=4$ Supersymmetric Gauge Theories and Monopoles of Spin 1}（Physics Letters B83, 1979, pp.~321--326）}
The central charges in the supersymmetry algebra of the $\cN=4$ supersymmetry gauge theory are obtained. When spontaneous symmetry breaking is imposed it is shown that the spins of the topological monopole states should be identical to those of the massive elementary particles.

\medskip
本文求出 $\cN=4$ 超对称规范理论中超对称代数的中心荷。在施加自发对称破缺之后，可以证明：拓扑单极子态的\emph{自旋}应当与有质量基本粒子的自旋相同。
\end{abstractbox}

\noindent $\cN=2$ 的情形后来怎样，我稍后会讲。

\subsection{唯一的确定胜利：$U(1)$ 与 Hodge 对偶}

\srcp{23}当时我对 Montonen--Olive 对偶仍然心存怀疑，但无论如何，有一点当时看来很清楚——而这一点我肯定是对的——即使这个对偶是正确的，用当时可用的方法也不可能在理解它方面取得重大进展。不过，在继续往下讲之前，我应当指出一个后来变得重要起来的基本事实。对于规范群为 $U(1)$ 的阿贝尔理论——确切地说，是没有带电场的“自由”理论——对偶实际上是正确的，而且在这种情形下不需要超对称。对 $U(1)$ 而言，Montonen--Olive 对偶是调和形式 Hodge 对偶的量子版本。换句话说，真空中的麦克斯韦方程 $\dd F=\dd\star F=0$ 说的正是：$U(1)$ 规范场的曲率 $F$ 是一个调和形式。这个条件在 Hodge 对偶
\[
  F\longrightarrow \star F
\]
之下不变；而对规范群 $U(1)$ 来说，Montonen--Olive 对偶正是这件事的量子版本。

\subsection{停滞的十年}

\srcp{24}但无论如何，此后许多年里，Montonen--Olive 对偶方面没有多少进展。在这段时期，Atiyah 正在——除别的事情之外——研究单极子。我想稍微解释一下那项工作。

\begin{pizhu}{中心荷、BPS 界与“经典质量公式”}
为什么“超对称”能救“经典公式”？关键在于超对称代数被拓扑荷修正。在 $\cN=2$ 理论里，代数的反对易子不只是动量，还多出中心荷 $Z$：
\[
  \{Q^i_{\ \dot A},\bar Q^{j}_{\ \dot B}\}=\varepsilon_{\dot A\dot B}\,\varepsilon^{ij}Z ,
\]
而对 BPS 态，质量恰好等于中心荷的模：$M=|Z|$。中心荷是拓扑量（单极子的磁荷、或电—磁组合），不随耦合常数跑；于是 $M=|Z|$ 这条等式在任意耦合下都成立：\emph{经典质量公式没有量子修正}。

这正是 Witten--Olive 的结果，也是它“不是对偶的证据”的原因：它只说明了质量谱的\emph{一部分}是刚性的，并没有说明两个理论（$G$ 与 $\gvee$）是同一个理论。用 Witten 后来的措辞：他们证明了 Bogomolny 界在量子层面成立并被饱和，而不是证明了对偶。要证对偶，还得把整个谱、以及算子代数都匹配上——这是第 6 章 Sen 的“smoking gun”才做到的事。
\end{pizhu}

\begin{pizhu}{$\cN=1,2,4$ 与三个除法代数}
原文说“极粗略地说，$\cN=1,2,4$ 对应复数、四元数、八元数”。这是一个助记，不是定理，值得知道它的边界在哪里。

四维 $\cN=1,2,4$ 的超荷数分别是 $4,8,16$；R-对称群分别是 $U(1)$、$U(2)$、$SU(4)$。把超对称与除法代数联系起来的严格版本出现在别的维数：在三维，$\cN=1,2,4,8$ 的超杨--米尔斯理论与 $\RR,\CC,\HH,\OO$ 相关；在四维，$16$ 个超荷的 $\cN=4$ 理论正是十维 $\cN=1$ 理论的维数约化。所以 Witten 用了“极粗略地”三个字：这句话给出的是“超对称越多，代数结构越像可除代数”的直觉，而不是一个可以逐项对照的表。

真正要紧的区别在下一句：$\cN=4$ 对给定规范群唯一，$\cN=2$ 不唯一（可以加物质、加超势），$\cN=1$ 更自由。这个“唯一性”是 $\cN=4$ 能成为对偶的干净试验场的根本原因——没有别的参数可以调，对偶要么对，要么错。
\end{pizhu}

\begin{pizhu}{$U(1)$ 的自由理论：唯一能算的对偶}
无电荷的 $U(1)$ 理论里，$F$ 是闭的且余闭的，故 $\star F$ 也是；把 $F$ 换成 $\star F$ 就是把电场与磁场互换。量子层面，这个操作把耦合 $e$ 送到 $4\pi/e$（在本书第 3 章的约定下），而自由理论的作用量在 $\tau\mapsto-1/\tau$ 下不变——这是可以逐项算出来的，不需要任何猜想。

Witten 后来反复用这个例子：它是唯一一个“电磁对偶可以被证明”的场论，也是所有关于对偶的直觉的校准点。它与 $\cN=4$ 的关系，正像调和形式与 Hodge 理论的关系：前者是后者的自由场特例；而 $\cN=4$ 的非阿贝尔对偶，则是把这个自由理论的对偶“非阿贝尔化”。
\end{pizhu}
